\begin{figure}[h]
    \centering
\begin{adjustbox}{max width=\textwidth, max height=0.92\textheight, keepaspectratio}
\begin{tikzpicture}[
    font=\small,
    >=Latex,
    node distance=0.7cm and 1.0cm,
    inputblock/.style={
        draw=black,
        fill=red!12,
        rounded corners=2pt,
        minimum width=2.8cm,
        minimum height=0.9cm,
        align=center
    },
    convblock/.style={
        draw=blue!70!black,
        fill=blue!15,
        rounded corners=2pt,
        minimum width=4.0cm,
        minimum height=1.0cm,
        align=center,
        thick
    },
    resblock/.style={
        draw=black!70,
        fill=gray!15,
        rounded corners=3pt,
        minimum width=4.3cm,
        minimum height=1.3cm,
        align=center,
        thick
    },
    opblock/.style={
        draw=black!70,
        fill=yellow!20,
        rounded corners=2pt,
        minimum width=3.5cm,
        minimum height=0.8cm,
        align=center
    },
    skipblock/.style={
        draw=black!70,
        fill=green!12,
        rounded corners=2pt,
        minimum width=3.5cm,
        minimum height=0.8cm,
        align=center
    },
    shapetext/.style={
        font=\scriptsize,
        align=center
    },
    mainarrow/.style={->, thick},
    skiparrow/.style={->, thick, dashed},
    projblock/.style={
        draw=purple!70!black,
        fill=purple!12,
        rounded corners=2pt,
        minimum width=3.7cm,
        minimum height=0.8cm,
        align=center
    },
    sidearrow/.style={->, thick, dashed}
]
% ============================================================
% MAIN ROW -- left-to-right data flow
% ============================================================
\node[inputblock] (input) {
    \textbf{Input}
};
% Moved slightly left by reducing distance from Input
\node[convblock, right=0.55cm of input] (inputproj) {
    \textbf{input\_projection}\\
    Conv1d, kernel size $1$
};
% Increased distance from input_projection so ResidualBlock remains visually separated
\node[resblock, right=2.75cm of inputproj] (residual) {
    \textbf{ResidualBlock}\\
    with skip taps
};
% Moved skip aggregation slightly to the right
\node[skipblock, right=1.75cm of residual] (sumskip) {
    $\displaystyle \sum \mathrm{skips}/\sqrt{n}$
};
\node[convblock, right=0.9cm of sumskip] (outputproj) {
    \textbf{output\_projection}\\
    Conv1d, kernel size $1$
};
\node[opblock, right=0.9cm of outputproj] (reshape) {
    $D_x = c_{\mathrm{skip}}\,x_t \;+\; c_{\mathrm{out}}\,F_x$
};
\node[inputblock, right=0.9cm of reshape] (output) {
    \textbf{Output}
};
% ============================================================
% DIMENSION LABELS -- upper-right corner of every block
% (input label removed per request)
% ============================================================
\node[shapetext, anchor=south west]
    at ($(inputproj.north east)+(-1.0,0.01)$)
    {$(B,\mathrm{ch},K,L)$};
\node[shapetext, anchor=south west]
    at ($(residual.north east)+(-1.0,0.01)$)
    {$(B,\mathrm{ch},K{\times}L)$};
\node[shapetext, anchor=south west]
    at ($(sumskip.north east)+(-1.0,0.01)$)
    {$(B,\mathrm{ch},K{\times}L)$};
\node[shapetext, anchor=south west]
    at ($(outputproj.north east)+(-1.5,0.01)$)
    {$(B,1,K{\times}L)$};
\node[shapetext, anchor=south west]
    at ($(reshape.north east)+(-1.0,0.01)$)
    {$(B,K,L)$};
\node[shapetext, anchor=south west]
    at ($(output.north east)+(-1.0,0.01)$)
    {$(B,K,L)$};
% ============================================================
% MAIN FLOW ARROWS -- left to right
% ============================================================
\draw[mainarrow] (input)      -- (inputproj);
\draw[mainarrow] (inputproj)  -- (residual);
\draw[mainarrow] (residual)   -- (sumskip);
\draw[mainarrow] (sumskip)    -- (outputproj);
\draw[mainarrow] (outputproj) -- (reshape);
\draw[mainarrow] (reshape)    -- (output);
% ============================================================
% EXTRA ARROWS FROM INPUT
% Arrow from top-center of input to top-center of ResidualBlock
% ============================================================
\draw[mainarrow]
    (input.north)
    -- ++(0, 1.4cm)
    -| (residual.north);
% Arrow from bottom-center of input to bottom-center of Reshape
\draw[mainarrow]
    (input.south)
    -- ++(0, -1.4cm)
    -| (reshape.south);
% ============================================================
% REPEATED-BLOCK ANNOTATION
% ============================================================
\node[anchor=west, font=\large]
    at ($(residual.east)+(10.0cm,1.5cm)$)
    {$\times n$};
% ============================================================
% DASHED CONTAINER -- enlarged to include input_projection,
% ResidualBlock, sumskip (skips), and output_projection.
% ============================================================
\node[
    draw=black,
    dashed,
    rounded corners=3.8pt,
    inner sep=0.55cm,
    fit=(inputproj)(residual)(sumskip)(outputproj)
] (rescontainer) {};
\end{tikzpicture}
\end{adjustbox} % end adjustbox
    \caption{High-level overview of the CSDI denoising network. The two-channel input (details in Fig.\ \ref{fig:model_input}) is projected into a latent space, processed by $n$ stacked residual blocks with skip connections, and decoded back to the predicted clean signal.}
    \label{fig:csdi_overview}
\end{figure}